\documentclass[11pt,a4paper]{article}

%% PACHETE MATEMATICE
\usepackage{amsmath,amssymb,amsfonts}
\usepackage{amsthm}
\usepackage{mathpazo}
\usepackage{eulervm}
\usepackage{enumerate}
\usepackage{color}
\usepackage{graphicx}
\usepackage[all]{xy}
\usepackage{xcolor}
\usepackage{framed}
\usepackage[unicode]{hyperref}
\setcounter{MaxMatrixCols}{20}
\usepackage{multicol}
%\usepackage[romanian]{babel}


%% ASPECT
\linespread{1.05}
\usepackage[T1]{fontenc}
\usepackage[utf8x]{inputenc}
\usepackage[margin=2cm]{geometry}
% \usepackage[color]{showkeys}
\usepackage{anyfontsize}
\usepackage{titlesec}
\titleformat*{\section}{\large\bfseries}

\usepackage{fancyhdr}
\pagestyle{fancy}
\rhead{\small \color{gray} Notițe de Adrian Manea}
\lhead{\small \color{gray} BCD-Aero, 2017---2018, L. Costache \& M. Olteanu}


\usepackage[autostyle = false, style = german]{csquotes}
\usepackage[romanian]{babel}
\newcommand\qq{\enquote}

%% COMENZILE MELE
\renewcommand*{\proofname}{Dem.:}
\newcommand\bb{\mathbb}
\newcommand\dr{\mathrm}
\newcommand\kal{\mathcal}
\newcommand\fk{\mathfrak}
\newcommand\op{\oplus}
\newcommand\ot{\otimes}
\newcommand\ds{\displaystyle}
\newcommand\id{\indent}
\newcommand\nid{\noindent}
\newcommand\seq{\subseteq}
\newcommand\sneq{\subsetneq}
\newcommand\speq{\supseteq}
\newcommand\spneq{\supsetneq}
\newcommand\rar{\rightarrow}
\newcommand\Rar{\Rightarrow}
\newcommand\xrar{\xrightarrow}
\newcommand\lar{\leftarrow}
\newcommand\Lar{\Leftarrow}
\newcommand\xlar{\xleftarrow}
\newcommand\lrar{\leftrightarrow}
\newcommand\Lrar{\Leftrightarrow}
\newcommand\ideal{\trianglelefteq}
\newcommand\ai{\text{ a.\^{i}. }}
\newcommand\sk{(\textit{Schi\c{t}\u{a}}) }
\newcommand\rhup{\rightharpoonup}
\newcommand\rhdn{\rightharpoondown}
\newcommand\lhup{\leftharpoonup}
\newcommand\lhdn{\leftharpoondown}
\newcommand\vid{\emptyset}
\newcommand\wh{\widehat}
\newcommand\ol{\overline}
\newcommand\NN{\mathbb{N}}
\newcommand\RR{\mathbb{R}}
\newcommand\ZZ{\mathbb{Z}}
\newcommand\QQ{\mathbb{Q}}
\newcommand\CC{\mathbb{C}}

%% STILURILE MELE
\newtheoremstyle{dotless}{}{}{\itshape}{}{\bfseries}{:}{ }{}

\theoremstyle{dotless}
\newtheorem{theorem}{Teorem\u{a}}[section]
\newtheorem{proposition}{Propozi\c{t}ie}[section]
\newtheorem{lemma}{Lem\u{a}}[section]
\newtheorem{corollary}{Corolar}[section]
\newtheorem{exercise}{Exerci\c{t}iu}

\newtheoremstyle{dotlessdr}{}{}{}{}{\bfseries}{:}{ }{}
\theoremstyle{dotlessdr}
\newtheorem{example}{Exemplu}[section]
\newtheorem{definition}{Defini\c{t}ie}[section]
\newtheorem{remark}{Observa\c{t}ie}[section]




\begin{document}

\begin{center}
  {\large\textbf{Seminar 13 \\ Recapitulare}}
\end{center}

\textbf{MODEL 1}

1. Să se calculeze cu o eroare mai mică decît $\varepsilon = 10^{-3}$:
\[
  \int_0^{\frac{1}{2}} \frac{1 - \cos x^2}{x^2} dx.
\]

\vspace{1cm}

2. Să se arate, folosind definiția, că funcția:
\[
  f(x, y) = \frac{2xy}{x^2 + y^2}, \quad (x, y) \neq (0, 0)
\]
nu are limită în origine.

\vspace{1cm}

3. Să se verifice dacă funcția $f(x, y) = (x^2 + y^2)^{-\frac{1}{2}}$ este
armonică.

\vspace{1cm}

4. Să se calculeze, folosind definiția, derivatele parțiale ale funcției $f$
în punctul $(-1, 1)$, unde:
\[
  f : \RR^2 \to \RR, \quad f(x, y) = 2x^3y - e^{x^2 + y^2}.
\]

\vspace{1cm}

5. Determinați valorile extreme pentru funcția $f$, definită pe domeniul
$D$, unde:
\[
  f(x, y) = x^4 + y^3 - 4x^3 - 3y^2 + 3y, \quad D = \{(x, y) \in \RR^2 \mid x^2 + y^2 \leq 4 \}.
\]

\vspace{2cm}

\textbf{MODEL 2}

1. Fie funcția $f(x) = \sqrt[4]{x + 1}$.
\begin{enumerate}[(a)]
\item Să se scrie polinoamele $T_1(x)$ și $T_2(x)$ în jurul punctului $a = 0$;
\item Să se calculeze $\sqrt[4]{10004}$ cu 4 zecimale exacte.
\end{enumerate}

\vspace{1cm}

2. Să se arate că funcția:
\[
  f(x, y) = \begin{cases}
    \frac{2xy}{x^2 + y^2}, & (x, y) \neq (0, 0) \\
    0, & (x, y) = (0, 0)
  \end{cases}
\]
este continuă în raport cu fiecare variabilă în parte, dar nu este continuă în raport
cu ansamblul variabilelor.

\vspace{1cm}

3. Fie $f(x, y) = \sqrt{\sin^2 x + \sin^2 y}$. Să se calculeze, folosind definiția,
$\frac{\partial f}{\partial x} (\frac{\pi}{4}, 0)$ și
$\frac{\partial f}{\partial y} (\frac{\pi}{4}, \frac{\pi}{4})$.

\vspace{1cm}

4. Fie seria de puteri:
\[
  \sum_{n \geq 0} (-1)^n \frac{x^{2n+1}}{2n + 1}.
\]

\begin{enumerate}[(a)]
\item Să se determine domeniul de convergență;
\item Să se determine suma seriei.
\end{enumerate}

\vspace{1cm}

5. Să se determine punctele de extrem și valorile extreme pentru
funcția
\[
  f : \RR^2 \to \RR, \quad f(x, y) = \sin x \sin y \sin(x + y).
\]

\vspace{2cm}

\textbf{MODEL 3}

1. Să se determine aproximările liniare, pătratice și cubice ale funcției
\[
  f(x) = x \ln x,
\]
în jurul punctului $a = 1$.

\vspace{1cm}

2. Fie funcția:
\[
  f(x, y) = \begin{cases}
    xy \sin \frac{x^2 - y^2}{x^2 + y^2}, & (x, y) \neq (0, 0) \\
    0, & (x, y) = (0, 0).
  \end{cases}
\]

Să se arate că $f \in \kal{C}^1(\RR^2)$.

\vspace{1cm}

3. Fie funcția:
\[
  f(x, y, z) = g(xy, xyz + \frac{x^2}{2} + \frac{x^4}{4}).
\]

Să se arate că ea satisface ecuația:
\[
  -xy \frac{\partial f}{\partial x} + y^2 \frac{\partial f}{\partial y} +
  x(1 + x^2) \frac{\partial f}{\partial x} = 0.
\]

\vspace{1cm}

4. Să se determine punctele cele mai apropiate de origine care se găsesc pe suprafața
de ecuație $4x^2 + y^2 + z^2 - 8x - 4y + 4 = 0$.

\vspace{1cm}

5. Fie seria de puteri $\ds\sum_{n \geq 1} \frac{(x - 1)^{2n}}{n \cdot 9^n}$.

Să se determine domeniul de convergență.

\vspace{1cm}

6. Să se determine suma seriei $\ds\sum_{n \geq 1} nx^{-n}$.

\newpage

\textbf{PARȚIAL SERIA A}

1. Fie seria de puteri $\ds\sum_{n \geq 1} \frac{1}{n} x^{2n}$.
Să se determine domeniul de convergență și suma seriei.

\vspace{1cm}

2. Să se calculeze cu o eroare de cel mult $\varepsilon = 10^{-2}$:
\[
  \int_0^{\frac{1}{4}} \frac{1 - e^{-x}}{x} dx.
\]

\vspace{1cm}

3. Fie funcția $f(x) = 1 + x \ln x$.
\begin{enumerate}[(a)]
\item Să se determine polinoamele Taylor $T_1, T_2$ în jurul lui $1$;
\item Aproximați $\ln \frac{11}{10}$, folosind $T_2$.
\end{enumerate}

\vspace{1cm}

4. Considerăm seria de puteri:
\[
  \sum_{n \geq 1} \frac{(-1)^{n-1}}{n} x^{2n}.
\]

\begin{enumerate}[(a)]
\item Găsiți domeniul de convergență.
\item Calculați suma seriei.
\end{enumerate}

\vspace{1cm}

5. Calculați valoarea integralei, cu o eroare mai mică decît $10^{-2}$:
\[
  \int_0^{\frac{1}{2}} \frac{1 - \cos x}{x^2} dx.
\]

\vspace{1cm}

6. Fie funcția $f(x) = (x + 1)e^x$.
\begin{enumerate}[(a)]
\item Scrieți $T_1$ și $T_2$ în jurul lui $0$;
\item Calculați cu aproximație $e^{-0,1}$, folosind expresia lui $T_2(x)$
  și găsiți un majorant pentru eroare.
\end{enumerate}

\end{document}
